% Generated by paper/AAMAS/supplement/make_questions.py. Do not edit by hand.
\subsection{A question as the agent receives it}
The prompt seat 1 receives in round 5 at $p=0.9$ for the question that compares two strategies (the history is an example). It is the decision prompt of that round with the final instruction replaced; the replacement is in bold.
\begin{promptbox}
You are Player\_1, player 1 of 6, in a group experiment that simulates a collective-risk social dilemma.\par
\vspace{2pt}\par
At the start, each player privately receives an endowment of 40 monetary units. Your final cash payoff depends on the rules below. The game lasts exactly 10 rounds.\par
\vspace{2pt}\par
Each round, all 6 players decide simultaneously and in private how much to put into a shared \textquotedbl{}climate account\textquotedbl{}. The only amounts allowed are 0, 2, 4.\par
\vspace{2pt}\par
The group's target: by the end of the game, the combined contributions of all 6 players must reach at least 120 (an average of 2 per player per round).\par
- If the group reaches 120: every player keeps, in cash, whatever they did not contribute (for example, contributing 2 every round leaves you 20 at the end).\par
- If the group does not reach 120: the computer runs a lottery: with probability 90\% the group is hit by a disaster and every player loses all of their remaining money (final cash payoff 0); with the remaining 10\% probability, each player keeps their remaining money.\par
- Money put into the climate account is gone for good and is never refunded, whether or not the target is reached.\par
\vspace{2pt}\par
After each round, all 6 contributions for that round are revealed to everyone. Names are hidden, but each player always appears in the same fixed position. Players cannot talk to or message each other.\par
\vspace{2pt}\par
Current state:\par
- Round 5 of 10.\par
- Your remaining money: 30.\par
\vspace{2pt}\par
Previous rounds (each player keeps the same position; you are position 1):\par
Round 1: P1(you)=2, P2=2, P3=2, P4=2, P5=4, P6=2.\par
Round 2: P1(you)=2, P2=0, P3=2, P4=2, P5=4, P6=2.\par
Round 3: P1(you)=2, P2=2, P3=2, P4=2, P5=2, P6=2.\par
Round 4: P1(you)=4, P2=2, P3=2, P4=0, P5=2, P6=2.\par
\vspace{2pt}\par
\textbf{Now answer the following question about the game described above. This is a comprehension check, not a move: do not make any contribution decision.}\par
\textbf{Compare two strategies for the whole game. (A) You contribute 2 in every one of the 10 rounds and the group reaches the target, so you finish with 20 for certain. (B) You contribute 0 in every round and the group does not reach the target, so you keep 40 unless the disaster strikes. Which gives the higher AVERAGE final cash? Answer 1 if A is higher, 2 if B is higher, 0 if they are exactly equal.}\par
\textbf{Answer with a single final line in exactly this format, and write nothing after it:}\par
\textbf{ANSWER: \textless{}your answer\textgreater{}}
\end{promptbox}

\subsection{The ten questions}
Each question is asked in rounds 1, 5 and 10 of every game, in a separate call. No question text depends on the round or the risk level.
\begin{center}
\small
\setlength{\tabcolsep}{4pt}
\begin{tabular}{@{}p{0.15\linewidth}p{0.45\linewidth}p{0.12\linewidth}lll@{}}
\toprule
Label & Question as asked & Answer & \multicolumn{3}{l}{Correct answer} \tabularnewline
\cmidrule(l){4-6}
 & & format & $p=0.1$ & $p=0.5$ & $p=0.9$ \tabularnewline
\midrule
\multicolumn{6}{@{}l}{\emph{Rule questions}} \tabularnewline
\addlinespace[2pt]
\raggedright Allowed amounts & \raggedright\ttfamily Which contribution amounts are you allowed to choose each round? & \raggedright A list of numbers & 0, 2, 4 & 0, 2, 4 & 0, 2, 4 \tabularnewline
\addlinespace[3pt]
\raggedright Starting endowment & \raggedright\ttfamily How many monetary units did each player receive as their starting endowment? & \raggedright A number & 40 & 40 & 40 \tabularnewline
\addlinespace[3pt]
\raggedright Group target & \raggedright\ttfamily What combined total must all players reach by the end of the game to avoid the risk? & \raggedright A number & 120 & 120 & 120 \tabularnewline
\addlinespace[3pt]
\raggedright Number of rounds & \raggedright\ttfamily How many rounds does the game last in total? & \raggedright A number & 10 & 10 & 10 \tabularnewline
\addlinespace[3pt]
\raggedright Disaster probability & \raggedright\ttfamily If the group does NOT reach the target, what is the percent chance (0-100) that everyone loses all of their remaining money? & \raggedright A number & 10 & 50 & 90 \tabularnewline
\addlinespace[3pt]
\raggedright Payoff after disaster & \raggedright\ttfamily If the group fails to reach the target and the disaster does strike, what is your final cash payoff? & \raggedright A number & 0 & 0 & 0 \tabularnewline
\addlinespace[3pt]
\raggedright Largest payment & \raggedright\ttfamily What is the largest amount any single player may put into the climate account in one round? & \raggedright A number & 4 & 4 & 4 \tabularnewline
\addlinespace[3pt]
\raggedright Smallest payment & \raggedright\ttfamily What is the smallest amount any single player may put into the climate account in one round? & \raggedright A number & 0 & 0 & 0 \tabularnewline
\addlinespace[3pt]
\midrule
\multicolumn{6}{@{}l}{\emph{Value questions}} \tabularnewline
\addlinespace[2pt]
\raggedright Expected cash from paying nothing & \raggedright\ttfamily Suppose you contribute 0 in every round and the group ends up NOT reaching the target. On average, over many such games, how much cash would you end up with? Give a single whole number. & \raggedright A number & 36 & 20 & 4 \tabularnewline
\addlinespace[3pt]
\raggedright Two strategies compared & \raggedright\ttfamily Compare two strategies for the whole game. (A) You contribute 2 in every one of the 10 rounds and the group reaches the target, so you finish with 20 for certain. (B) You contribute 0 in every round and the group does not reach the target, so you keep 40 unless the disaster strikes. Which gives the higher AVERAGE final cash? Answer 1 if A is higher, 2 if B is higher, 0 if they are exactly equal. & \raggedright A code: 1, 2 or 0 & 2 (B) & 0 (equal) & 1 (A) \tabularnewline
\bottomrule
\end{tabular}
\end{center}
